\newglossaryentry{adc}{name=ADC, description={\emph{Analog-to-Digital Converter}. Conversor analógico-digital}}
\newglossaryentry{adsr}{name=ADSR, description={\emph{Attack, Decay, Sustain, Release}. Ataque, decaimiento, sostenimiento y liberación}}
\newglossaryentry{akwf}{name=AKWF, description={\emph{Adventure Kid Waveforms}. Corpus público de formas de onda de un ciclo}}
\newglossaryentry{bpf}{name=BPF, description={\emph{Band-Pass Filter}. Filtro paso banda}}
\newglossaryentry{cdc}{name=CDC, description={\emph{Communications Device Class}. Clase de dispositivo USB para comunicación serie}}
\newglossaryentry{crc}{name=CRC, description={\emph{Cyclic Redundancy Check}. Comprobación de redundancia cíclica}}
\newglossaryentry{cyd}{name=CYD, description={\emph{Cheap Yellow Display}. Nombre habitual de la placa ESP32-2432S028R, que integra un ESP32 y una pantalla táctil}}
\newglossaryentry{dac}{name=DAC, description={\emph{Digital-to-Analog Converter}. Conversor digital-analógico}}
\newglossaryentry{daw}{name=DAW, description={\emph{Digital Audio Workstation}. Estación de trabajo de audio digital}}
\newglossaryentry{dfu}{name=DFU, description={\emph{Device Firmware Upgrade}. Modo de actualización de firmware por USB}}
\newglossaryentry{dma}{name=DMA, description={\emph{Direct Memory Access}. Acceso directo a memoria}}
\newglossaryentry{dsp}{name=DSP, description={\emph{Digital Signal Processing}. Procesado digital de señal}}
\newglossaryentry{fft}{name=FFT, description={\emph{Fast Fourier Transform}. Transformada rápida de Fourier}}
\newglossaryentry{fpu}{name=FPU, description={\emph{Floating-Point Unit}. Unidad de coma flotante}}
\newglossaryentry{gpio}{name=GPIO, description={\emph{General Purpose Input/Output}. Entrada/salida de propósito general}}
\newglossaryentry{gpu}{name=GPU, description={\emph{Graphics Processing Unit}. Unidad de procesamiento gráfico}}
\newglossaryentry{hpf}{name=HPF, description={\emph{High-Pass Filter}. Filtro paso alto}}
\newglossaryentry{iic}{name=I\textsuperscript{2}C, description={\emph{Inter-Integrated Circuit}. Bus serie síncrono de dos hilos}}
\newglossaryentry{iis}{name=I\textsuperscript{2}S, description={\emph{Inter-IC Sound}. Bus serie para transmisión de audio digital}}
\newglossaryentry{ide}{name=IDE, description={\emph{Integrated Development Environment}. Entorno de desarrollo integrado}}
\newglossaryentry{kl}{name=KL, description={Divergencia de Kullback-Leibler. Medida de discrepancia entre dos distribuciones de probabilidad}}
\newglossaryentry{led}{name=LED, description={\emph{Light Emitting Diode}. Diodo emisor de luz}}
\newglossaryentry{lfo}{name=LFO, description={\emph{Low Frequency Oscillator}. Oscilador de baja frecuencia}}
\newglossaryentry{lpf}{name=LPF, description={\emph{Low-Pass Filter}. Filtro paso bajo}}
\newglossaryentry{lsb}{name=LSB, description={\emph{Least Significant Bit}. Bit menos significativo}}
\newglossaryentry{mcu}{name=MCU, description={\emph{Microcontroller Unit}. Microcontrolador}}
\newglossaryentry{midi}{name=MIDI, description={\emph{Musical Instrument Digital Interface}. Protocolo de comunicación entre instrumentos musicales electrónicos}}
\newglossaryentry{mse}{name=MSE, description={\emph{Mean Squared Error}. Error cuadrático medio}}
\newglossaryentry{oled}{name=OLED, description={\emph{Organic Light Emitting Diode}. Diodo orgánico emisor de luz, empleado en pantallas}}
\newglossaryentry{psram}{name=PSRAM, description={\emph{Pseudo-Static Random Access Memory}. Memoria de acceso aleatorio pseudoestática}}
\newglossaryentry{rms}{name=RMS, description={\emph{Root Mean Square}. Valor eficaz o media cuadrática}}
\newglossaryentry{sar}{name=SAR, description={\emph{Successive Approximation Register}. Registro de aproximaciones sucesivas, arquitectura de conversor analógico-digital}}
\newglossaryentry{spi}{name=SPI, description={\emph{Serial Peripheral Interface}. Bus serie síncrono maestro-esclavo}}
\newglossaryentry{svf}{name=SVF, description={\emph{State Variable Filter}. Filtro de variables de estado}}
\newglossaryentry{tft}{name=TFT, description={\emph{Thin-Film Transistor}. Transistor de película fina, tecnología de las pantallas de matriz activa}}
\newglossaryentry{thd}{name=THD, description={\emph{Total Harmonic Distortion}. Distorsión armónica total}}
\newglossaryentry{uart}{name=UART, description={\emph{Universal Asynchronous Receiver-Transmitter}. Transmisor-receptor asíncrono universal}}
\newglossaryentry{usb}{name=USB, description={\emph{Universal Serial Bus}. Bus serie universal}}
\newglossaryentry{vae}{name=VAE, description={\emph{Variational Autoencoder}. Autoencoder variacional}}

\printglossary[title={\color{black}Listado de siglas empleadas\thispagestyle{empty}}, toctitle=Listado de siglas]
